% ==========================================================================
% Preamble for the 2026-10-21 Emerging Technologies Taskforce talk.
% A copy of the INFO 4617 Web Data Science course preamble
% (cuinfoscience/INFO4617-Fall2026, slides/common/preamble.tex), so this
% talk looks like the course decks. Two changes from the course copy:
%   - no biblatex: the talk cites sources in \bottomcite lines, so the
%     build needs no biber;
%   - no \cupids macro, which needs the course's logo file.
% The Gotham theme files beside this one are copied unchanged from the
% same folder.
% ==========================================================================

\usetheme{gotham}

\usepackage{libertinus}
\usepackage{booktabs}
\usepackage{graphicx}
\usepackage{hyperref}
% Resolve images whether a deck is compiled from its own folder (img/...)
% or from the repo root, as Overleaf does (slides/week-NN/img/...). Both
% dialects appear in instructor- and agent-authored frames; make both build.
% This talk sits one folder below the repo root, so ../ resolves a path
% written from the root (2026-10-21-ett/img/...), as Overleaf's file tree shows it.
\graphicspath{{./}{../}{../../}}
\usepackage{tikz}
\usepackage{xcolor}
\usepackage{multirow}
\usepackage{listings}
\usetikzlibrary{positioning, arrows.meta, shapes}
\usepackage{hwemoji}
\DeclareUnicodeCharacter{FE0F}{}


% Named vertical-space vocabulary for use in frame bodies, in place of
% ad hoc \vspace{Nem} calls. LaTeX's kernel defaults for \smallskip /
% \medskip / \bigskip (3pt/6pt/12pt) are tuned for 10-12pt body text and
% land too large at this deck's 9pt base size and \tiny/\footnotesize/
% \small frame content -- \bigskip alone would exceed every \vspace value
% previously used in this repo. Recalibrated here to match the deck's
% established rhythm instead.
\setlength{\smallskipamount}{3pt plus 1pt minus 1pt}
\setlength{\medskipamount}{6pt plus 1.5pt minus 1.5pt}
\setlength{\bigskipamount}{12pt plus 2pt minus 2pt}

\gothamset{sectiontocframe default=off}

% ------------------------------------------------------------------ colors
\definecolor{cugold}{HTML}{CFB87C}
\definecolor{tab-red}{HTML}{E15759}
\definecolor{tab-orange}{HTML}{F28E2B}
\definecolor{tab-blue}{HTML}{4E79A7}
\definecolor{tab-cyan}{HTML}{76B7B2}
\definecolor{tab-green}{HTML}{59A14F}
\definecolor{tab-yellow}{HTML}{EDC948}
\definecolor{tab-purple}{HTML}{B07AA1}
\definecolor{tab-pink}{HTML}{FF9DA7}
\definecolor{tab-brown}{HTML}{9C755F}
\definecolor{tab-gray}{HTML}{BAB0AC}

\setbeamercolor{frametitle}{fg=black, bg=cugold}
\setbeamercolor{progress bar}{fg=cugold, bg=cugold!20}

\setbeamercolor{block title alerted}{fg=black, bg=cugold}
\setbeamercolor{block body alerted}{fg=black, bg=cugold!10}

% ------------------------------------------------------------------ footline
% This talk's copy differs from the course decks' here. The page number and
% \bottomcite share one baseline, \footdepth above the slide's bottom edge,
% and both sit inside the text margins: the citation starts where body text
% starts (\beamer@leftmargin), and the page number ends where body text ends
% (\beamer@rightmargin). The footer box has fixed heights in points, so the
% citation's baseline can be set to the same length.
\makeatletter
\newlength{\footdepth}\setlength{\footdepth}{8pt}
% What the footer's right end says; the appendix renews it to "Appendix"
\newcommand{\footlabel}{\insertframenumber{} of \inserttotalframenumber}
\setbeamertemplate{footline}{%
  \begin{beamercolorbox}[wd=\paperwidth,ht=9pt,dp=\footdepth,right]{footline}%
    {\footnotesize\footlabel}\hspace*{\beamer@rightmargin}%
  \end{beamercolorbox}%
}

% ------------------------------------------------------------------ helpers
% \bottomcite{text} -- a small attribution or citation, set on the page
% number's baseline at the left text margin, with a thin rule above it. A
% TikZ overlay anchored to the physical page places it, so it doesn't depend
% on how much content precedes it and takes no room from the frame body. It
% is a bottom-aligned \parbox, so a citation that wraps grows upward and its
% last line stays on the page number's baseline; its width stops 5em short of
% the right text margin, leaving room for the page number.
\newcommand{\bottomcite}[1]{%
  \begin{tikzpicture}[remember picture, overlay]
    \node[anchor=base west, inner sep=0pt]
      at ([xshift=\beamer@leftmargin, yshift=\footdepth]current page.south west)
      {\parbox[b]{\dimexpr\paperwidth-\beamer@leftmargin-\beamer@rightmargin-5em\relax}{%
        \raggedright\tiny\rule{\linewidth}{0.3pt}\\#1}};
  \end{tikzpicture}%
}
\makeatother

% \flushitems -- put first-level bullets flush with the surrounding text, so a
% list sits directly under its heading with no indent. Beamer draws each bullet
% right-aligned, ending \labelsep before \leftmargini, so a \leftmargini of the
% bullet's own width plus \labelsep starts the bullet at the text's left edge.
% Use it inside a group (a column, a frame) before the itemize it should affect.
\newcommand{\flushitems}{%
  \settowidth{\leftmargini}{\usebeamerfont*{itemize item}\usebeamertemplate{itemize item}}%
  \addtolength{\leftmargini}{\labelsep}%
}

% Consistent code styling for the Wednesday "applied notebook" frames.
\lstdefinestyle{wds}{
    basicstyle=\ttfamily\scriptsize,
    keywordstyle=\color{tab-blue}\bfseries,
    commentstyle=\color{tab-green},
    stringstyle=\color{tab-orange},
    showstringspaces=false,
    breaklines=true,
    columns=fullflexible,
    language=Python,
    aboveskip=0.4em,
    belowskip=0.4em,
    xleftmargin=0pt,
    xrightmargin=0pt
}
\lstset{style=wds}

% \dayband{Day}{Topic} -- a lightweight section divider label used to mark
% Monday / Wednesday / Friday within a week's deck.
\newcommand{\dayband}[2]{{\large\textbf{#1}}\;\textcolor{tab-gray}{\textbar}\;#2}
